\section{From construction order to optimization geometry}
\label{sec:dynamics}

The scalar loss order is a visibility order: a nonzero leading term need not provide descent on the admissible transition cone. Let
\[
L(x)-L(0)=P_d(x)+P_{d+1}(x)+\cdots
\]
with homogeneous $P_d\not\equiv0$, and let $K$ be the cone of admissible local perturbation rays. Define $\dvis=d$ and let $\ddown$ be the least first-visible degree along a ray $v\in K$ whose leading coefficient is negative. Then
\begin{proposition}
\label{prop:descent}
$\ddown\geq\dvis$, with equality iff $P_{\dvis}(v)<0$ for some admissible ray $v$.
\end{proposition}
This follows directly from the one-dimensional expansion $L(tv)-L(0)=t^{d(v)}P_{d(v)}(v)+o(t^{d(v)})$. Our bootstrap fixtures satisfy the sign condition explicitly: their isolated leading term is
\[
L-L_0=-C\prod_{i=1}^d x_i,\qquad C>0,
\]
on the positive construction cone, so $\ddown=\dloss=d$.

\paragraph{Affine probability coordinates.}
For independent coordinates initialized symmetrically at $x_i(0)=\delta>0$, Euclidean gradient flow preserves $x_i=x$ and gives $\dot x=Cx^{d-1}$. The exact time to a fixed threshold $\theta>\delta$ is
\[
\tau_d(\delta,\theta)=
\begin{cases}
(\theta-\delta)/C,&d=1,\\
C^{-1}\log(\theta/\delta),&d=2,\\
\big(\delta^{2-d}-\theta^{2-d}\big)/(C(d-2)),&d\geq3.
\end{cases}
\]
Thus the small-initialization classes are finite, logarithmic, and $\Theta(\delta^{-(d-2)})$. These homogeneous-flow exponents are known dynamics; the finite-memory result tells us which degree $d$ a dormant topology exposes.

A fixed positive diagonal preconditioner changes the balanced geometry and prefactor but not this degree-controlled class. More generally, a smooth local reparameterization $\varepsilon=\phi(\theta)$ with nonsingular Jacobian $J=D\phi(0)$ preserves scalar vanishing order because $P_d(\phi(\theta))=P_d(J\theta)+O(\|\theta\|^{d+1})$ and $P_d\circ J$ is nonzero. This does \emph{not} make Euclidean trajectories coordinate invariant.

\paragraph{Simplex boundary.}
The absent-edge point is singular for a softmax/logit chart. For the same isolated degree-$d$ construction with $p_i=\sigma(z_i)$ and symmetric rare-edge probability $p$, Euclidean logit flow gives
\[
\dot p=Cp^{d+1}(1-p)^2,
\]
so for fixed positive threshold
\[
\boxed{\tau_d^{\mathrm{logit}}(\delta)=\frac{1}{Cd}\delta^{-d}(1+o(1)).}
\]
For a full $K$-way softmax target edge,
\[
\|\nabla_z p\|^2=p^2\left((1-p)^2+\sum_aq_a^2\right),
\]
and
\[
\frac{K}{K-1}p^2(1-p)^2\leq\|\nabla_zp\|^2\leq2p^2(1-p)^2.
\]
Hence an isolated symmetric degree-$d$ target construction still has $\dot p=\Theta(p^{d+1})$ near the rare-edge boundary for fixed $K$. The absolute bootstrap exponent is therefore parameterization-dependent, but a one-unit scaffold reduction in $d$ removes one rare-edge power from the asymptotic time in this geometry as well.

\begin{figure}[t]
\centering
\maybegraphics[width=\linewidth]{figures/figure4_order_to_time_geometry.pdf}
\caption{Construction order fixes local homogeneity, not an optimizer-invariant clock. The same degree maps to different completion-time classes under affine probability and rare-edge softmax Euclidean flow.}
\label{fig:geometry}
\end{figure}

The intended causal chain is
\[
\boxed{\text{dormant topology}\to\text{construction order}\to\text{optimizer geometry}\to\text{bootstrap time}.}
\]
Appendix~\ref{app:dynamics} gives the exact integrations, repeated-parameter extension, and singular-chart derivations.
